\chapter{Digital Twin in Healthcare and eHealth}

\section{Digital Twin for Health}
A Digital Twin for Health inherits the general physical--virtual structure introduced in the Digital Twin fundamentals chapter, but it also depends on the eHealth ecosystem described in the previous chapter. Healthcare twins add domain-specific requirements that are particularly important in clinical contexts: they must preserve clinical meaning, integrate multimodal patient data, support longitudinal updating, communicate uncertainty, protect privacy, and remain subject to validation and human supervision. A Digital Twin for Health is therefore a virtual representation used to support healthcare-related monitoring, simulation, prediction, and decision-making while remaining connected to a clearly identified patient, organ, device, process, or health system.

Depending on its scope, it may represent a patient, an organ, a disease process, a medical device, a clinical workflow, a hospital unit, or a public health system. Unlike a simple dashboard, a Digital Twin aims to maintain a dynamic relationship between real-world data and virtual models. The general Digital Twin criteria still apply: a domain-specific health twin should not be considered mature merely because it contains medical data or an AI model.

\section{Types of Healthcare Digital Twins}
Healthcare Digital Twins can be classified into several categories:

\begin{itemize}[leftmargin=*]
    \item \textbf{Patient-specific Digital Twins:} virtual representations of individual patients using multimodal data.
    \item \textbf{Organ or body-system Digital Twins:} models focusing on a specific organ or physiological system.
    \item \textbf{Disease Digital Twins:} models representing disease progression, risk factors, or treatment response.
    \item \textbf{Medical device Digital Twins:} virtual representations of devices used for design, monitoring, or safety assessment.
    \item \textbf{Hospital and process Digital Twins:} models of workflows, resources, patient flow, logistics, and healthcare units.
    \item \textbf{Public health Digital Twins:} models used for population-level monitoring, planning, prevention, and crisis management.
\end{itemize}

\begin{figure}[H]
    \centering
    \resizebox{0.95\textwidth}{!}{\begin{tikzpicture}[font=\small, >=Latex, node distance=1.0cm]
\tikzstyle{root}=[rounded corners=8pt, draw=DTBlue, very thick, fill=DTLightBlue, minimum width=5.3cm, minimum height=1.1cm, align=center]
\tikzstyle{leaf}=[rounded corners=6pt, draw=DTGray, thick, fill=DTLightGray, minimum width=3.6cm, minimum height=0.95cm, align=center]
\tikzstyle{arrow}=[->, thick, DTGray]

\node[root] (root) {\textbf{Digital Twins in Healthcare and eHealth}};
\node[leaf, below left=1.2cm and 3.2cm of root] (patient) {Patient-specific\\digital twins};
\node[leaf, below=1.2cm of root] (organ) {Organ / body-system\\digital twins};
\node[leaf, below right=1.2cm and 3.2cm of root] (disease) {Disease and care-pathway\\digital twins};
\node[leaf, below=1.4cm of patient] (hospital) {Hospital and workflow\\digital twins};
\node[leaf, below=1.4cm of organ] (device) {Medical device\\digital twins};
\node[leaf, below=1.4cm of disease] (public) {Public health and\\population twins};

\foreach \x in {patient,organ,disease,hospital,device,public}
  \draw[arrow] (root) -- (\x);
\end{tikzpicture}
}
    \caption{General taxonomy of Digital Twin applications in healthcare and eHealth. Author synthesis based on the reviewed literature.}
    \label{fig:healthcare_dt_taxonomy}
\end{figure}

\section{Applications in eHealth}
The applications of Digital Twins in eHealth are diverse. They include remote monitoring, personalized care pathways, treatment planning, rehabilitation support, chronic disease management, hospital resource optimization, medical device development, drug discovery, clinical trial simulation, and public health planning. These applications share a common objective: to transform heterogeneous data into actionable knowledge.

\section{Patient-Centered eHealth Digital Twins}
A patient-centered eHealth Digital Twin integrates clinical, physiological, behavioral, and contextual data to support personalized monitoring and decision-making. Such a system may receive information from EHRs, connected devices, laboratory results, patient-reported outcomes, and environmental sources. However, the level of clinical reliability depends on data quality, model validation, and the ability to handle uncertainty.

\section[Non-Invasive Identification of Patient-Specific Digital Twins]{Non-Invasive Identification of Patient-Specific\\ Digital Twins}
Many clinically relevant physiological variables are difficult to measure directly. Some variables require invasive procedures, specialized imaging, or controlled clinical settings, while routine eHealth systems usually observe only indirect signals such as biosignals, images, wearable measurements, laboratory values, or clinician-entered summaries. A central challenge for patient-specific Digital Twins is therefore to infer a plausible internal state from partial and non-invasive observations without confusing inference with measurement.

In this context, Artificial Intelligence can support an inverse-identification task. Instead of directly mapping a medical observation to a clinical label, the learning system estimates latent patient-specific parameters that initialize or calibrate a mechanistic model. The mechanistic model then simulates physiological variables such as pressure, flow, volume, or state trajectories. This differs from direct medical prediction: a model that maps echocardiography to ejection fraction or a disease class may be useful, but it does not necessarily create a patient-specific simulator. Digital Twin identification follows a richer chain: non-invasive observation, estimated physiological parameters, mechanistic model, and simulated patient state.

Med-Real2Sim, introduced by Kuang et al.~\cite{kuang2024medreal2sim}, illustrates this emerging direction in cardiovascular modeling. The method uses non-invasive echocardiography videos to infer parameters of a mechanistic cardiovascular model and then reconstruct pressure--volume behavior that would otherwise require more invasive characterization. Its importance for healthcare Digital Twins is methodological: AI does not replace physiology, but helps identify patient-specific parameters under physiological constraints.

Such approaches remain limited by identifiability and validation. A simulated state may match an observable value such as ventricular volume while still depending on non-unique or model-dependent latent parameters. For this reason, non-invasive Digital Twin identification must preserve the distinction between measured observations, estimated parameters, simulated states, and predicted or counterfactual outputs.

\section{Healthcare Process Digital Twins}
Digital Twins are not limited to individual patients. They can also represent healthcare processes, such as emergency department flow, outpatient services, operating room scheduling, hospital bed allocation, or vaccination campaigns. In this context, the Digital Twin supports simulation of operational scenarios and optimization of healthcare resources.

\section{Ethical, Legal, and Social Considerations}
Healthcare Digital Twins use sensitive data and may influence medical decisions. Therefore, their development must consider consent, privacy, data ownership, transparency, fairness, bias mitigation, accountability, and clinical responsibility. Ethical design is not an optional component; it is a necessary condition for trustworthy eHealth Digital Twins.

\section{Challenges and Limitations}
The implementation of Digital Twins in eHealth faces several challenges: interoperability between healthcare systems, access to high-quality multimodal data, real-time synchronization, computational cost, cybersecurity, explainability, regulatory approval, integration into clinical workflows, and acceptance by healthcare professionals and patients.

\section{Summary}
Digital Twins offer a promising direction for eHealth, but their adoption requires more than technical models. It requires architecture, governance, validation, security, interoperability, ethical alignment, and careful handling of non-invasive patient-specific model identification. These issues are addressed in the state-of-the-art chapters that follow.

\clearpage
